% ------------------------------------------------------------------------------
% 8 Results V -- Heterogeneous treatment effects
% ------------------------------------------------------------------------------
\section{Results V: Heterogeneous Treatment Effects}\label{sec:res-cate}

\subsection{Average effects on the five-year absolute risk}

Table~\ref{tab:cate} reports the average treatment effect of levamisole
plus 5-FU versus observation on the five-year DFS event risk, estimated
three ways on the jackknife pseudo-observations ($n = 594$ complete
cases). The three estimators agree closely: the unadjusted
Kaplan--Meier difference is $-16.8$ percentage points (bootstrap 95\%
CI $-24.5$ to $-9.1$), cross-fitted G-computation gives $-15.3$
points, and the doubly robust AIPW estimator gives $-16.2$ points with
influence-function standard error $3.9$ (95\% CI $-23.9$ to $-8.5$).
The consistency across a design-free estimator, an outcome-model-based
estimator, and a doubly robust estimator is exactly what one hopes to
see in a randomized comparison, and the magnitudes match the ITT
analysis of Section~\ref{sec:res-tx}.

\begin{table}[htbp]
\centering\small
\caption{Average treatment effect on the five-year DFS event risk
(negative $=$ risk reduction with levamisole plus 5-FU).}
\label{tab:cate}
\begin{tabular*}{\textwidth}{@{\extracolsep{\fill}} l c c c}
\toprule
Estimator & ATE (percentage points) & SE & 95\% CI \\
\midrule
Kaplan--Meier difference (bootstrap)   & $-16.8$ & 4.0 & $-24.5$ to $-9.1$ \\
G-computation (cross-fitted, pseudo-obs) & $-15.3$ & 0.5 & $-16.2$ to $-14.4$ \\
AIPW (cross-fitted, doubly robust)      & $\mathbf{-16.2}$ & 3.9 & $-23.9$ to $-8.5$ \\
\bottomrule
\end{tabular*}
\end{table}

\subsection{Conditional effects: where the benefit lives}

The cross-fitted T-learner decomposes this average into
patient-level predicted benefits
(Figure~\ref{fig:catedist}). The distribution is wide: predicted
five-year risk reductions range from negligible to over 40 percentage
points, with quartile means of $28.9$, $18.6$, $12.6$, and $1.1$
points. Crucially, the predicted benefit is strongly tied to baseline
risk (correlation $0.648$ between control-arm predicted risk and
predicted benefit; right panel of Figure~\ref{fig:catedist}): the
patients predicted to benefit most are those whose disease is most
aggressive---high nodal burden, advanced local extent---while
low-risk patients have little room for absolute improvement even
under a proportionally constant treatment effect.

\begin{figure}[htbp]
  \centering
  \includegraphics[width=0.92\textwidth, max height=0.38\textheight]{fig_cate_distribution.png}
  \caption{Predicted absolute benefit of levamisole plus 5-FU at five
  years. Left: distribution across patients (AIPW average marked).
  Right: predicted benefit versus predicted baseline risk, coloured by
  randomized arm.}
  \label{fig:catedist}
\end{figure}

Table~\ref{tab:cateq} sharpens the picture with observed outcomes:
within the highest predicted-benefit quartile (mean baseline risk
$71.0\%$), the observed five-year risk difference is $-12.8$ points;
within the lowest quartile (baseline risk $46.6\%$), it is $-19.3$
points---statistically compatible with the same relative effect
everywhere. The apparent paradox (large predicted but modest observed
benefit in Q1, the reverse in Q4) reflects both estimation noise in
quartile-specific contrasts (each based on $\sim$64--98 events) and
the mechanical relationship between baseline risk and achievable
absolute benefit. The Kaplan--Meier curves within the extreme
quartiles (Figure~\ref{fig:catekm}) visualize the trade-off: in the
high-benefit quartile both arms fall steeply but the curves separate;
in the low-benefit quartile both arms do comparatively well.

\begin{table}[htbp]
\centering\footnotesize
\caption{Predicted-benefit quartiles from the cross-fitted T-learner.
``Predicted benefit'' is the model-implied five-year risk reduction;
``observed risk difference'' is the corresponding pseudo-observation
contrast by arm within quartile.}
\label{tab:cateq}
\begin{tabular*}{\textwidth}{@{\extracolsep{\fill}} l c c c c c}
\toprule
Quartile & $n$ & Pred.\ benefit (pts) & Baseline risk & Observed\ diff (pts) & Events \\
\midrule
Q1 (largest predicted benefit) & 149 & 28.9 & 71.0\% & $-12.8$ & 98 \\
Q2                             & 148 & 18.6 & 58.2\% & $-16.1$ & 79 \\
Q3                             & 148 & 12.6 & 52.0\% & $-11.8$ & 72 \\
Q4 (smallest)                  & 149 & 1.1  & 46.6\% & $-19.3$ & 64 \\
\bottomrule
\end{tabular*}
\end{table}

\begin{figure}[htbp]
  \centering
  \includegraphics[width=0.92\textwidth, max height=0.38\textheight]{fig_cate_km_quartiles.png}
  \caption{Kaplan--Meier disease-free survival within the lowest and
  highest predicted-benefit quartiles.}
  \label{fig:catekm}
\end{figure}

The synthesis is the section's key message: in this trial, heterogeneity
in \emph{absolute} benefit is real and large, but it is driven almost
entirely by heterogeneity in \emph{baseline risk} rather than by
qualitative effect modification---consistent with the null interaction
tests of Section~\ref{sec:res-design}. For clinical decision-making,
this is the most useful kind of heterogeneity: it means a single
relative effect, applied to an accurate prognostic model, suffices to
identify who gains most in absolute terms, with no need to invoke
distinct treatment-response biology. All such CATE statements remain
exploratory in a single trial of this size, and the honest framing is
that the analysis demonstrates the \emph{methods} and generates
hypotheses for external validation.
